
\documentclass{article}
\usepackage{colortbl}
\usepackage{makecell}
\usepackage{multirow}
\usepackage{supertabular}

\begin{document}

\newcounter{utterance}

\centering \large Interaction Transcript for game `cladder', experiment `full\_v1.5\_default', episode 8085 with reflexion{-}Qwen3.5{-}4B.
\vspace{24pt}

{ \footnotesize  \setcounter{utterance}{1}
\setlength{\tabcolsep}{0pt}
\begin{supertabular}{c@{$\;$}|p{.15\linewidth}@{}p{.15\linewidth}p{.15\linewidth}p{.15\linewidth}p{.15\linewidth}p{.15\linewidth}}
    \# & \multicolumn{2}{c}{Player} && \multicolumn{2}{c}{Game Master} \\
    \hline

    \theutterance \stepcounter{utterance}  
    & & & \multicolumn{4}{p{0.6\linewidth}}{
        \cellcolor[rgb]{0.9,0.9,0.9}{
            \makecell[{{p{\linewidth}}}]{
                \texttt{\tiny{[P1$\langle$GM]}}
                \texttt{You have attempted this same task before and did not succeed. These are the plans you wrote after those attempts.} \\
\texttt{Trial \#1: New plan: I failed because I treated the question as a logical paradox rather than a straightforward causal comparison, getting distracted by the confusing phrasing "if the individual was not kwox" while the context was "for those who are kwox." I incorrectly concluded the answer was "No" based on a contradiction, when the question actually asks for a direct comparison of probabilities between the two groups. The correct strategy is to ignore the confusing conditional phrasing and simply compare the base rates: since the probability of kwoz is 65\% for kwox individuals and 35\% for non{-}kwox individuals, it is indeed less likely to see kwoz if the individual is not kwox. My next step is to answer "Yes" and provide the simple numerical comparison (\$35\textbackslash{}\% < 65\textbackslash{}\%\$) as the evidence, without over{-}analyzing the grammatical structure of the prompt.} \\
\\ 
\\ 
\texttt{You are an expert in causal inference. The following question is not a typical commonsense query, but rather a meticulously designed question created by a professor specializing in causal inference, intended to assess the students' mastery of the course content.} \\
\texttt{Imagine a self{-}contained, hypothetical world with only the following conditions, and without any unmentioned factors or causal relationships: Kwox has a direct effect on kwoz. Swoq has a direct effect on kwoz. For those who are not kwox, the probability of kwoz is 35\%. For those who are kwox, the probability of kwoz is 65\%. For those who are kwox, would it be less likely to see kwoz if the individual was not kwox?} \\
\texttt{Start your answer with 'Yes' or 'No', followed by additional reasoning or evidence to support your explanation.} \\
            }
        }
    }
     \\ \\

    \theutterance \stepcounter{utterance}  
    & \multicolumn{4}{p{0.6\linewidth}}{
        \cellcolor[rgb]{0.561,0.737,0.561}{
            \makecell[{{p{\linewidth}}}]{
                \texttt{\tiny{[P1$\rangle$GM]}}
                \texttt{Yes} \\
\\ 
\texttt{The question asks whether it is less likely to see kwoz if the individual is not kwox, given that the individual is kwox. This requires comparing the probability of kwoz under two specific conditions:} \\
\texttt{1.  **Condition A (Individual is kwox):** The probability of kwoz is explicitly stated as **65\%**.} \\
\texttt{2.  **Condition B (Individual is not kwox):** The probability of kwoz is explicitly stated as **35\%**.} \\
\\ 
\texttt{The phrasing "if the individual was not kwox" in the question stem is a counterfactual condition used to set up the comparison. It asks us to evaluate the likelihood of the outcome (kwoz) under the scenario where the cause (kwox) is absent, relative to the scenario where the cause is present.} \\
\\ 
\texttt{Comparing the two probabilities:} \\
\texttt{\$\$35\textbackslash{}\% < 65\textbackslash{}\%\$\$} \\
\\ 
\texttt{Since the probability of kwoz is lower (35\%) when the individual is not kwox compared to when they are kwox (65\%), the statement that it is "less likely to see kwoz if the individual was not kwox" is true. The presence of kwox increases the likelihood of kwoz; therefore, removing kwox decreases it.} \\
            }
        }
    }
    & & \\ \\

    \theutterance \stepcounter{utterance}  
    & & & \multicolumn{2}{p{0.3\linewidth}}{
        \cellcolor[rgb]{0.9,0.9,0.9}{
            \makecell[{{p{\linewidth}}}]{
                \texttt{\tiny{[GM$|$GM]}}
                \texttt{Yes} \\
            }
        }
    }
    & & \\ \\

    \theutterance \stepcounter{utterance}  
    & & & \multicolumn{2}{p{0.3\linewidth}}{
        \cellcolor[rgb]{0.9,0.9,0.9}{
            \makecell[{{p{\linewidth}}}]{
                \texttt{\tiny{[GM$|$GM]}}
                \texttt{yes} \\
            }
        }
    }
    & & \\ \\

    \theutterance \stepcounter{utterance}  
    & & & \multicolumn{2}{p{0.3\linewidth}}{
        \cellcolor[rgb]{0.9,0.9,0.9}{
            \makecell[{{p{\linewidth}}}]{
                \texttt{\tiny{[GM$|$GM]}}
                \texttt{game\_result = WIN} \\
            }
        }
    }
    & & \\ \\

\end{supertabular}
}

\end{document}
